% ============================================
% MEM-26
% Gerado de molambudos.md — build_miolo.py
% ============================================
\section*{\textbf{MEM-26} --- A Última Página}
\markboth{MEM-26 --- A Última Página}{}

\noindent \textit{3 dias antes de morrer}

\noindent \textbf{1979 --- Três dias antes de morrer}

\vspace{0.5em}\noindent\rule{\textwidth}{0.5pt}\vspace{0.5em}

\noindent Eu sabia que ia morrer.

\noindent O corpo estava parando. O coração falhava. A respiração era curta. A criatura estava pesada demais dentro de mim.

\noindent Três dias antes de morrer, escrevi a última página.

\noindent \textit{\textquotedblleft{}Quem ler isto: você é o próximo.}

\noindent \textit{Eu carreguei a criatura por 62 anos. Ela entrou em mim quando eu era criança, numa vala, no meio de um campo de concentração. Ela me comeu por dentro. Mas ela também me manteve vivo.}

\noindent \textit{Ela precisa de hospedeiro. Quando eu morrer, ela vai procurar outro.}

\noindent \textit{Não um corpo. Um leitor.}

\noindent \textit{Você vai ler este diário. Enquanto lê, ela vai passar para você. Não fisicamente. Espiritualmente. A história vai entrar na sua cabeça. E, uma vez dentro, não sai mais.}

\noindent \textit{Você vai começar a sonhar com a vala. Vai sentir o cheiro. Vai ver o olho amarelo no canto do quarto.}

\noindent \textit{Vai achar que é sugestão.}

\noindent \textit{Não é.}

\noindent \textit{É a criatura.}

\noindent \textit{Bem-vindo ao ciclo, meu substituto.}

\noindent \textit{--- Joaquim Antônio Correia}
\textit{--- Paciente 1.260}
\textit{--- Última palavra: fome.\textquotedblright{}}\NE{O editor registra: a página original deste fragmento foi encontrada solta no bolso interno do paletó de Joaquim, durante a exumação de 1979. O papel estava coberto de terra e umidade; a tinta havia borrado em diversos pontos, especialmente na margem inferior. A transcrição aqui reproduzida é o texto legível após restauro. A última palavra --- \textquotedblleft{}fome\textquotedblright{} --- estava quase ilegível, preservada apenas pela pressão da caneta sobre o papel, que marcou a página seguinte mesmo após a tinta se dissolver.}

\vspace{0.5em}
\noindent\textit{\textbf{↪ Rotas:} \rota{DOC-06} \rotasep \rota{DOC-08}}
